\chapter{Related Work}
%
\label{chap:RelatedWork}
%
This chapter surveys the research areas most directly relevant to this thesis:
(1)~the use of LLMs for information extraction, (2)~automated service ticket
analysis, (3)~multi-label text classification, and (4)~multilingual NLP.

%----------------------------------------------------------------------
\section{Large Language Models for Information Extraction}
\label{sec:rw-llm-ie}
%
Since the introduction of the Transformer architecture~\cite{vaswani2017attention},
pre-trained language models have become the dominant paradigm for NLP tasks.
\textsc{BERT}~\cite{devlin2019bert} demonstrated that fine-tuning a single
pre-trained model achieves strong results across diverse extraction tasks.
More recently, instruction-tuned decoder-only models such as the
GPT~\cite{brown2020language} and Llama~\cite{touvron2023llama} families have shown
that well-crafted zero-shot prompts can match or surpass fine-tuned baselines on
many information extraction benchmarks.

\citeauthor{wei2022finetuned}~\cite{wei2022finetuned} showed that instruction
fine-tuning substantially improves zero-shot performance, while
\citeauthor{kojima2022large}~\cite{kojima2022large} demonstrated that chain-of-thought
prompting further unlocks reasoning in larger models.  Both findings motivate the
structured, step-by-step prompt design used in this thesis
(see Section~\ref{sec:prompt-design}).

Domain-aware prompts consistently outperform generic prompts.
\citeauthor{ma2023large}~\cite{ma2023large} showed this for biomedical named entity
recognition, and similar results have been reported in legal document analysis.
Accordingly, the prompts in this thesis embed Amazone-specific domain vocabulary and
define a fixed output schema.

%----------------------------------------------------------------------
\section{Automated Service Ticket Analysis}
\label{sec:rw-tickets}
%
Automated analysis of IT helpdesk tickets has been studied since at least the
early 2000s.  Classical approaches relied on keyword matching and TF-IDF
vectorisation combined with support vector machines for ticket categorisation
\cite{yang1999reevaluation}.  More recent work applies deep learning: recurrent
networks \cite{hochreiter1997long} and BERT-based classifiers are now common for
ticket routing and priority prediction.

However, the bulk of this literature targets English-language IT helpdesk data
from publicly available datasets~(e.g.~Incident Ticket from ServiceNow).
Industrial service tickets for physical machinery—multi-lingual, attachment-heavy,
and spanning years of email threads—are comparatively understudied.
The closest related work by \citeauthor{hassan2023llm}~\cite{hassan2023llm}
applies GPT-4 to manufacturing quality reports but does not address the
extraction of structured entities or multi-label categorisation.

This thesis contributes to filling this gap by constructing and releasing a
benchmark for agricultural machinery service tickets and evaluating multiple LLMs
on it.

%----------------------------------------------------------------------
\section{Multi-label Text Classification}
\label{sec:rw-multilabel}
%
Service tickets require multi-label classification because a single ticket may
simultaneously involve a warranty claim, a spare part request, and an electronic
fault.  Classical multi-label learning surveys~\cite{zhang2014review} distinguish
problem-transformation methods~(e.g.\ binary relevance, label powerset) from
algorithm-adaptation methods.

Transformer models reformulate multi-label classification as sequence generation,
outputting a comma-separated list of labels or a JSON array.  The blind-tag
evaluation protocol used in this thesis follows the approach of
\citeauthor{yang2023llm4label}~\cite{yang2023llm4label}, who found that
LLMs outperform fine-tuned classifiers when the label set is large and the
training data is scarce—precisely the setting of this thesis.

An important distinction is between \emph{accuracy-oriented} evaluation~(strict
match against a fixed ground truth) and \emph{usefulness-oriented} evaluation.
The latter is more appropriate when multiple valid labellings exist, and is
captured here through the LLM-as-judge protocol described in
Section~\ref{sec:llm-judge}.

%----------------------------------------------------------------------
\section{Multilingual NLP and Code-switching}
\label{sec:rw-multilingual}
%
Service tickets in this corpus are predominantly in German but frequently contain
English phrases, Latin abbreviations, and brand-specific terminology.
This phenomenon, known as \emph{code-switching}, poses challenges for models
trained predominantly on English data.

Multilingual models such as mBERT~\cite{devlin2019bert} and
XLM-R~\cite{conneau2020unsupervised} handle code-switching with varying success
depending on the language pair and domain.  Among the LLMs evaluated in this
thesis, instruction-tuned models with multilingual training data~(Gemma~3,
Qwen~2.5) generally outperform English-centric models on German passages
(see Section~\ref{sec:results-language}).

%----------------------------------------------------------------------
\section{LLM Evaluation and the LLM-as-Judge Framework}
\label{sec:rw-judge}
%
Standard NLP metrics such as F1-score and exact match are insufficient for
evaluating free-form generation where multiple valid outputs exist.
\citeauthor{zheng2024judging}~\cite{zheng2024judging} introduced the
LLM-as-judge paradigm, in which a capable model~(e.g.~GPT-4) evaluates the
quality of outputs from other models using structured rubrics.  This approach
has been adopted for summarisation, dialogue, and, more recently, information
extraction evaluation.

This thesis applies LLM-as-judge to rate tag usefulness and reasoning coherence,
complementing the F1-based evaluation against human annotations.

%----------------------------------------------------------------------
\section{Summary and Research Gap}
\label{sec:rw-gap}
%
Table~\ref{tab:rw-gap} summarises the positioning of this thesis relative to
prior work.

\begin{table}[htbp]
    \centering
    \caption{Positioning of this thesis in the related-work landscape.}
    \label{tab:rw-gap}
    \begin{tabular}{lccccc}
        \hline
        \textbf{Work} & \textbf{LLM} & \textbf{Industrial} & \textbf{Multi-lingual}
                      & \textbf{Multi-label} & \textbf{Benchmark} \\
        \hline
        Yang~et~al.~\cite{yang1999reevaluation}   & -- & -- & -- & \checkmark & \checkmark \\
        Hassan~et~al.~\cite{hassan2023llm}         & \checkmark & \checkmark & -- & -- & -- \\
        Yang~et~al.~\cite{yang2023llm4label}       & \checkmark & -- & -- & \checkmark & \checkmark \\
        \textbf{This thesis}                        & \checkmark & \checkmark & \checkmark & \checkmark & \checkmark \\
        \hline
    \end{tabular}
\end{table}

The combination of an industrial, multi-lingual setting with multi-label
classification and a human-annotated benchmark represents a contribution not
addressed by any single prior work.
